\chapter{Release 4 --- Monitoring et supervision}

\section*{Introduction}
\addcontentsline{toc}{section}{Introduction}
Cette dernière release met en place la supervision des microservices déployés sur le cluster Kubernetes. L'objectif est de collecter les métriques des services, de les visualiser à travers des tableaux de bord et de configurer des alertes en cas d'anomalie.

%======================================================================
% SPRINT 7
%======================================================================
\section{Sprint 7 --- Prometheus, Grafana et dashboards}

\subsection{Objectif du sprint}
Déployer une stack de monitoring complète sur le cluster EKS pour superviser la santé et les performances des microservices en production.

\subsection{Sprint Backlog}

\begin{table}[H]
\centering
\renewcommand{\arraystretch}{1.2}
\begin{tabular}{|c|p{9cm}|c|c|}
\hline
\textbf{ID} & \textbf{Tâche} & \textbf{Priorité} & \textbf{Points} \\
\hline
7.1 & Déploiement de Prometheus sur le cluster Kubernetes & Should & 3 \\
\hline
7.2 & Déploiement de Grafana et configuration des sources de données & Should & 3 \\
\hline
7.3 & Création des dashboards de supervision des microservices & Should & 5 \\
\hline
7.4 & Configuration des ServiceMonitors pour le scraping des métriques & Should & 3 \\
\hline
\end{tabular}
\caption{Sprint Backlog --- Sprint 7}
\label{tab:sprint7_backlog}
\end{table}

\subsection{Architecture de monitoring}
La stack de monitoring est composée de trois outils complémentaires :
\begin{itemize}
    \item \textbf{Prometheus :} Système de collecte de métriques. Il interroge (scrape) périodiquement les endpoints \texttt{/actuator/prometheus} exposés par chaque microservice Spring Boot.
    \item \textbf{Grafana :} Plateforme de visualisation. Elle se connecte à Prometheus comme source de données et permet de créer des tableaux de bord interactifs.
    \item \textbf{Alertmanager :} Gestionnaire d'alertes. Il reçoit les alertes de Prometheus et les achemine vers les canaux de notification configurés.
\end{itemize}

\begin{figure}[H]
    \centering
    \fbox{\parbox{0.85\textwidth}{\centering\vspace{3cm}\textbf{[Diagramme d'architecture de monitoring]}\\\textit{Microservices (/actuator/prometheus) $\rightarrow$ Prometheus $\rightarrow$ Grafana}\vspace{3cm}}}
    \caption{Architecture de la stack de monitoring}
    \label{fig:monitoring_architecture}
\end{figure}

\subsection{Déploiement de Prometheus}
Prometheus est déployé sur le cluster Kubernetes dans un namespace dédié \texttt{monitoring}. Les \texttt{ServiceMonitor} (ressources personnalisées) sont configurés pour cibler automatiquement les microservices portant le label \texttt{app.kubernetes.io/part-of: billing-platform}.

\subsection{Dashboards Grafana}
Plusieurs tableaux de bord ont été créés pour visualiser les métriques clés :
\begin{itemize}
    \item \textbf{Vue globale :} Nombre de pods actifs, utilisation CPU et mémoire par service.
    \item \textbf{Requêtes HTTP :} Taux de requêtes, temps de réponse (p50, p95, p99) et taux d'erreurs par microservice.
    \item \textbf{JVM :} Utilisation du heap, nombre de threads actifs et activité du garbage collector.
\end{itemize}

\begin{figure}[H]
    \centering
    \fbox{\parbox{0.85\textwidth}{\centering\vspace{2.5cm}\textbf{[Capture d'écran --- Dashboard Grafana]}\\\textit{Vue globale des microservices}\vspace{2.5cm}}}
    \caption{Dashboard Grafana de supervision}
    \label{fig:grafana_dashboard}
\end{figure}

\begin{figure}[H]
    \centering
    \fbox{\parbox{0.85\textwidth}{\centering\vspace{2.5cm}\textbf{[Capture d'écran --- Prometheus Targets]}\\\textit{Liste des endpoints scrapés}\vspace{2.5cm}}}
    \caption{Targets Prometheus}
    \label{fig:prometheus_targets}
\end{figure}

\section*{Conclusion}
\addcontentsline{toc}{section}{Conclusion}
La mise en place de la stack de monitoring complète cette dernière release. L'ensemble de la plateforme de facturation télécom est désormais opérationnelle en production, avec une supervision continue des performance et de la santé des microservices.

\newpage
